\chapter{Lineare Gleichungssysteme}
\label{chap:linear_systems}

In vielen Anwendungen müssen mehrere lineare Bedingungen gleichzeitig erfüllt
sein. Im Maschinenbau bilden lineare Gleichungssysteme (kurz: LGS) die
Grundlage zahlreicher Berechnungsverfahren,
etwa in der Statik, bei der Finite-Elemente-Methode und innerhalb numerischer
Crashsimulationen. Weitere Anwendungen finden sich bei der Berechnung elektrischer
Netzwerke, von Strömungsmodellen und bei der Ausgleichsrechnung aus Messdaten.


\subsection*{Motivation}

Eine Last mit der Gewichtskraft $F_G=\qty{100}{\newton}$ wird von zwei Seilen
gehalten. Das erste Seil trägt $\qty{20}{\newton}$ mehr als das zweite. 
Um die unbekannten Seilkräfte $F_1$ und $F_2$ zu bestimmen, formulieren wir zwei
lineare Bedingungen:

\noindent
\begin{minipage}[c]{0.28\linewidth}
  \centering
  \CachedTikzFigure{content/resources/chapters/linear_systems/tikz/force_balance_mass.tex}
\end{minipage}%
\hfill
\begin{minipage}[c]{0.69\linewidth}
  \begin{ClozeArea}[24]{10}
    \centering
    \CachedTikzFigure{content/resources/chapters/linear_systems/tikz/force_balance_equations.tex}
  \end{ClozeArea}
\end{minipage}

Beide linearen Gleichungen müssen gleichzeitig erfüllt sein. Wie man leicht sehen kann, ist das der Fall für
\[
  F_1=\Cloze{2cm}{60}\,\unit{\newton}
  \quad \text{und} \quad
  F_2=\Cloze{2cm}{40}\,\unit{\newton}.
\]

\par\vspace{1cm}
In diesem Kapitel wollen wir die folgenden offenen Fragen klären:
\begin{itemize}
  \item Wie kann man systematisch die Lösungen eines LGS bestimmen?
  \item Wie viele Lösungen hat ein LGS?
\end{itemize}

\clearpage
\section{Grundbegriffe}
\label{sec:linear_systems_fundamentals}

Zunächst klären wir, was lineare Gleichungen und lineare Gleichungssysteme
sind und wann ein Tupel von Zahlen ein solches System löst.

\subsection{Lineare Gleichungen und lineare Gleichungssysteme}
\label{subsec:linear_systems_motivation_and_definition}

\begin{DefinitionBox}{Lineare Gleichung}{linear_systems_linear_equation}
  Eine \emph{lineare Gleichung} \EnglishTerm{linear equation} in den
  Unbekannten $x_1,\dots,x_n$ besitzt die Form
  \[
    a_1x_1+a_2x_2+\dots+a_nx_n=b.
  \]
  Die Zahlen $a_1,\dots,a_n$ heißen \emph{Koeffizienten}
  \EnglishTerm{coefficients}, $b$ heißt \emph{Absolutglied}
  \EnglishTerm{constant term}.
\end{DefinitionBox}

Die Unbekannten werden in einer linearen Gleichung weder miteinander
multipliziert noch potenziert und stehen nicht in nichtlinearen Funktionen.
Konstanten dürfen dagegen beliebig dargestellt sein:

\begin{ExampleBox}{Linear oder nichtlinear?}{linear_systems_linearity_examples}
  \begin{enumerate}[label=\alph*),itemsep=0.6em]
    \item $2x_1+3x_2=2^2$ ist \Cloze{3.2cm}{linear}.
    \item $x_1x_2=1$ ist \Cloze{3.2cm}{nicht linear}.
    \item $x_1^2+\sin(x_2)=1$ ist \Cloze{3.2cm}{nicht linear}.
  \end{enumerate}
\end{ExampleBox}

Fasst man mehrere lineare Gleichungen zusammen, so spricht man von einem
linearen Gleichungssystem.
\begin{DefinitionBox}{Lineares Gleichungssystem}{linear_systems_system}
  Ein \emph{lineares Gleichungssystem} \EnglishTerm{system of linear equations},
  kurz LGS, aus $m$ Gleichungen in $n$ Unbekannten besitzt die Form
  \[
    \begin{aligned}
      a_{11}x_1+a_{12}x_2+\dots+a_{1n}x_n&=b_1,\\
      a_{21}x_1+a_{22}x_2+\dots+a_{2n}x_n&=b_2,\\
      &\mathrel{\phantom{=}}\vdots\\
      a_{m1}x_1+a_{m2}x_2+\dots+a_{mn}x_n&=b_m.
    \end{aligned}
  \]
  Man spricht kurz von einem $m\times n$-LGS.
\end{DefinitionBox}

\begin{DefinitionBox}{Lösung und Lösungsmenge}{linear_systems_solution_set}
  Ein Tupel $(x_1,\dots,x_n)\in\R^n$ heißt \emph{Lösung}
  \EnglishTerm{solution} eines LGS, wenn es alle Gleichungen gleichzeitig
  erfüllt. Die Menge aller Lösungen heißt \emph{Lösungsmenge}
  \EnglishTerm{solution set} $\mathcal{L}$.
\end{DefinitionBox}


\clearpage
\section{Das Gauß-Verfahren}
\label{sec:linear_systems_gaussian_elimination}

Nun untersuchen wir Umformungen, welche die Lösungsmenge eines LGS nicht
verändern. Damit wollen wir ein Gleichungssystem schrittweise in eine Form
bringen, aus der sich seine Lösungen leicht ablesen lassen, ohne die Lösungen
dabei zu verändern.

\subsection{Erweiterte Koeffizientenmatrix}
\label{subsec:linear_systems_augmented_matrix}

Die Koeffizienten und Absolutglieder eines LGS lassen sich in einem rechteckigen
Zahlenschema anordnen. Die ersten $n$ Spalten bilden die
\emph{Koeffizientenmatrix} \EnglishTerm{coefficient matrix} $A$, die letzte
Spalte enthält den Vektor $b$ der Absolutglieder. Zusammen bilden sie die
\emph{erweiterte Koeffizientenmatrix} \EnglishTerm{augmented matrix}
$[A\mid b]$.

\begin{center}
  \CachedTikzFigure{content/resources/chapters/linear_systems/tikz/augmented_matrix_anatomy.tex}
\end{center}

Diese Notation verzichtet auf eine explizite Nennung der Unbekannten $x_j$.
Ihre Anzahl entspricht der Anzahl der Spalten von $A$.

Die Symbole $A$ und $b$ dienen hier ausschließlich als kompakte Notation;
Matrizen werden später ausführlich behandelt.

\begin{ExampleBox}{Vom LGS zur erweiterten Matrix}{linear_systems_augmented_matrix_example}
  Wir überführen die LGS in ihre erweiterte Koeffizientenmatrix $[A\mid b]$.
  \begin{itemize}[leftmargin=1.4em,itemsep=0.8em]
    \item
      \noindent
      \begin{minipage}[c]{0.47\linewidth}
        \[
          \begin{aligned}
             x_1-x_2       &=1,\\
            2x_1+3x_2+4x_3&=5
          \end{aligned}
        \]
      \end{minipage}%
      \begin{minipage}[c]{0.08\linewidth}
        \centering $\Longleftrightarrow$
      \end{minipage}%
      \begin{minipage}[c]{0.42\linewidth}
        \begin{ClozeArea}[12]{4}
          \centering
          \[
            \left[
              \begin{array}{@{}rrr|r@{}}
                1&-1&0&1\\
                2& 3&4&5
              \end{array}
            \right]
          \]
        \end{ClozeArea}
      \end{minipage}

    \item
      \noindent
      \begin{minipage}[c]{0.47\linewidth}
        \[
          \begin{aligned}
             2x_1+x_2&= 0,\\
              -x_1+3x_2&= 7,\\
                   4x_2&=-2
          \end{aligned}
        \]
      \end{minipage}%
      \begin{minipage}[c]{0.08\linewidth}
        \centering $\Longleftrightarrow$
      \end{minipage}%
      \begin{minipage}[c]{0.42\linewidth}
        \begin{ClozeArea}[12]{6}
          \centering
          \[
            \left[
              \begin{array}{@{}rr|r@{}}
                 2&1& 0\\
                -1&3& 7\\
                 0&4&-2
              \end{array}
            \right]
          \]
        \end{ClozeArea}
      \end{minipage}
  \end{itemize}
\end{ExampleBox}

\subsection{Elementare Zeilenumformungen}
\label{subsec:linear_systems_elementary_row_operations}

Die folgenden Umformungen ändern die Lösungsmenge nicht:
\[
  \begin{alignedat}{2}
    Z_i&\leftrightarrow Z_j
      &\qquad&\text{zwei Zeilen vertauschen},\\[0.45em]
    Z_i&\leftarrow \lambda Z_i,\quad \lambda\neq0
      &&\text{eine Zeile skalieren},\\[0.45em]
    Z_i&\leftarrow Z_i+\lambda Z_j
      &&\text{ein Vielfaches einer anderen Zeile addieren}.
  \end{alignedat}
\]

Sie heißen \emph{elementare Zeilenumformungen} \EnglishTerm{elementary row
operations}. Das entstehende LGS ist zum Ausgangssystem äquivalent und besitzt
dieselbe Lösungsmenge.

\begin{ExampleBox}{Die drei elementaren Zeilenumformungen}{linear_systems_eliminate_variable}
  Das folgende Beispiel enthält alle drei elementaren Zeilenumformungen. Die
  jeweils rechts angegebene Operation wird auf die nebenstehende Matrix
  angewendet. Die Lösungsmenge bleibt dabei unverändert.
  \[
    \begin{array}{r@{\qquad}l@{\qquad}l}
      \left[\begin{array}{@{}rr|r@{}}
        1&1&2\\
        0&1&1
      \end{array}\right]
      &\TransformOp{\Cloze{2.0cm}{Z_1\leftrightarrow Z_2}}
      &\mathcal L=\{(x_1,x_2)\in\R^2\mid x_1=1\land x_2=1\}\\[1.2em]
      \left[\begin{array}{@{}rr|r@{}}
        \Cloze{0.55cm}{0}&\Cloze{0.55cm}{1}&\Cloze{0.55cm}{1}\\
        \Cloze{0.55cm}{1}&\Cloze{0.55cm}{1}&\Cloze{0.55cm}{2}
      \end{array}\right]
      &\begin{array}{@{}l@{}}
        \smash{\TransformOp{\Cloze{1.1cm}{2Z_1}}}\\
        \smash{\TransformOp{\Cloze{1.1cm}{-3Z_2}}}
      \end{array}
      &\mathcal L=\{(x_1,x_2)\in\R^2\mid x_1=1\land x_2=1\}\\[1.2em]
      \left[\begin{array}{@{}rr|r@{}}
        \Cloze{0.55cm}{0}&\Cloze{0.55cm}{2}&\Cloze{0.55cm}{2}\\
        \Cloze{0.55cm}{-3}&\Cloze{0.55cm}{-3}&\Cloze{0.55cm}{-6}
      \end{array}\right]
      &\begin{array}{@{}l@{}}
        {}\\
        \smash{\TransformOp{\Cloze{2.0cm}{Z_2+Z_1}}}
      \end{array}
      &\mathcal L=\{(x_1,x_2)\in\R^2\mid x_1=1\land x_2=1\}\\[1.2em]
      \left[\begin{array}{@{}rr|r@{}}
        \Cloze{0.55cm}{0}&\Cloze{0.55cm}{2}&\Cloze{0.55cm}{2}\\
        \Cloze{0.55cm}{-3}&\Cloze{0.55cm}{-1}&\Cloze{0.55cm}{-4}
      \end{array}\right]
      &
      &\mathcal L=\{(x_1,x_2)\in\R^2\mid x_1=1\land x_2=1\}
    \end{array}
  \]
\end{ExampleBox}

\subsection{Gauß-Verfahren und Zeilenstufenform}
\label{subsec:linear_systems_gaussian_algorithm}

Die Grundidee des \emph{Gauß-Verfahrens} \EnglishTerm{Gaussian elimination}
besteht darin, die erweiterte Matrix durch elementare Zeilenumformungen in
Zeilenstufenform zu bringen. Aus dieser Form lassen sich die Unbekannten leicht
bestimmen.

\begin{DefinitionBox}{Zeilenstufenform}{linear_systems_row_echelon_form}
  Eine Matrix besitzt \emph{Zeilenstufenform} \EnglishTerm{row echelon form}
  genau dann, wenn
  \begin{enumerate}[label=\arabic*.,itemsep=0.25em]
    \item alle \emph{Nullzeilen} \EnglishTerm{zero rows} unterhalb der
      \emph{Nichtnullzeilen} \EnglishTerm{nonzero rows} stehen,
    \item der erste von null verschiedene Eintrag jeder tieferen Zeile weiter
      rechts steht als derjenige der vorherigen Zeile und
    \item unter diesen Einträgen nur Nullen stehen.
  \end{enumerate}

  \begin{center}
    \CachedTikzFigure{content/resources/chapters/linear_systems/tikz/row_echelon_staircase.tex}
  \end{center}
  Dabei kennzeichnet $\bullet$ den ersten von null verschiedenen Eintrag einer
  Nichtnullzeile; $\ast$ steht für einen beliebigen Eintrag, der auch null sein
  darf.
\end{DefinitionBox}


\begin{TheoremBox}{Gauß-Verfahren}{linear_systems_gaussian_algorithm}
  Das Gauß-Verfahren dient dazu, ein LGS systematisch in eine Form zu
  überführen, aus der sich seine Lösung leicht ablesen lässt. Es besteht aus
  zwei Schritten:
  \begin{enumerate}[label=\arabic*.,itemsep=0.3em]
    \item Bei der \emph{Vorwärtselimination} \EnglishTerm{forward elimination}
      wird die erweiterte Koeffizientenmatrix $[A\mid b]$ mithilfe elementarer
      Zeilenumformungen in Zeilenstufenform überführt.
    \item Beim \emph{Rückwärtseinsetzen} \EnglishTerm{back substitution}
      werden die Unbekannten ausgehend von der letzten Gleichung schrittweise
      bestimmt.
  \end{enumerate}
  \begin{center}
    \CachedTikzFigure{content/resources/chapters/linear_systems/tikz/gaussian_elimination_flow.tex}
  \end{center}
\end{TheoremBox}

\noindent\begin{minipage}{\linewidth}
\begin{ExampleBox}{Gauß-Verfahren}{linear_systems_gaussian_example}
  Wir lösen das folgende LGS mit dem Gauß-Verfahren und überführen es zunächst
  in die erweiterte Koeffizientenmatrix:
    \[
      \begin{aligned}
        2x_1+x_2-x_3&=3,\\
        3x_1+5x_2-4x_3&=1,\\
        4x_1-3x_2+2x_3&=2
      \end{aligned}
      \qquad\Longleftrightarrow\qquad
      \left[
        \begin{array}{@{}rrr|r@{}}
          \Cloze{0.7cm}{2}&\Cloze{0.7cm}{1}&\Cloze{0.7cm}{-1}&\Cloze{0.7cm}{3}\\[0.35em]
          \Cloze{0.7cm}{3}&\Cloze{0.7cm}{5}&\Cloze{0.7cm}{-4}&\Cloze{0.7cm}{1}\\[0.35em]
          \Cloze{0.7cm}{4}&\Cloze{0.7cm}{-3}&\Cloze{0.7cm}{2}&\Cloze{0.7cm}{2}
        \end{array}
      \right].
    \]
  \begin{ClozeArea}[34]{36}
    \textbf{Vorwärtselimination:} Die Operationen stehen jeweils rechts neben
    der Matrix, auf die sie angewendet werden:
    \[
      \begin{array}{r@{\;}r@{\qquad}l}
        &\left[\begin{array}{@{}rrr|r@{}}
          2& 1&-1& 3\\
          3& 5&-4& 1\\
          4&-3& 2& 2
        \end{array}\right]
        &\begin{array}{@{}l@{}}
          {}\\
          \smash{\TransformOp{2Z_2-3Z_1}}\\
          \smash{\TransformOp{Z_3-2Z_1}}
        \end{array}\\[1.0em]
        \Longleftrightarrow
        &\left[\begin{array}{@{}rrr|r@{}}
          2&1&-1& 3\\
          \textcolor{THMBlue}{0}&7&-5&-7\\
          \textcolor{THMBlue}{0}&-5&4&-4
        \end{array}\right]
        &\begin{array}{@{}l@{}}
          {}\\
          {}\\
          \smash{\TransformOp{7Z_3+5Z_2}}
        \end{array}\\[1.0em]
        \Longleftrightarrow
        &\left[\begin{array}{@{}rrr|r@{}}
          2&1&-1&3\\
          \textcolor{THMBlue}{0}&7&-5&-7\\
          \textcolor{THMBlue}{0}&\textcolor{THMBlue}{0}&3&-63
        \end{array}\right]
        &\text{Zeilenstufenform}
      \end{array}
    \]
    \textbf{Rückwärtseinsetzen:} Die Unbekannten werden von der letzten
    Gleichung aus nach oben bestimmt:
    \begin{itemize}[leftmargin=1.4em,itemsep=0.55em,topsep=0.3em]
      \item in Zeile 3 einsetzen:\qquad
        \(
          3x_3=-63
          \;\Longleftrightarrow\;
          x_3=\frac{-63}{3}
          \;\Longleftrightarrow\;
          x_3=-21
        \)
      \item in Zeile 2 einsetzen: \qquad
        \(
          7x_2-5x_3=-7
          \;\overset{x_3=-21}{\Longleftrightarrow}\;
          7x_2-5(-21)=-7
          \;\Longleftrightarrow\;
          x_2=-16
        \)
      \item in Zeile 1 einsetzen: \qquad
        \(
          2x_1+x_2-x_3=3
          \;\overset{\substack{x_2=-16\\x_3=-21}}{\Longleftrightarrow}\;
          2x_1+(-16)-(-21)=3
          \;\Longleftrightarrow\;
          x_1=-1
        \)
    \end{itemize}
    Somit lautet die eindeutige Lösung
    \[
      \mathcal{L}=\{(-1,-16,-21)\}.
    \]
  \end{ClozeArea}
\end{ExampleBox}

In diesem Beispiel hatte das LGS eine Lösung. Muss das immer der Fall sein?
\end{minipage}


\clearpage
\section{Lösungstypen und Lösungsmengen}
\label{sec:linear_systems_solution_sets}

Das Gauß-Verfahren hat im vorigen Abschnitt genau eine Lösung geliefert.
Zunächst untersuchen wir geometrisch, welche anderen Lösungsmengen auftreten
können. Anschließend erkennen wir den jeweiligen Fall algebraisch an der
Zeilenstufenform.

\subsection{Drei geometrische Lösungstypen}
\label{subsec:linear_systems_solution_sets_and_geometry}

Die Lösungsmenge eines LGS kann sehr unterschiedlich aussehen. Wir
untersuchen die möglichen Fälle zunächst an $2\times2$-LGS.

In zwei Dimensionen beschreibt jede lineare Gleichung eine Gerade in der
$x_1$-$x_2$-Ebene. Ein Lösungstupel $(x_1,x_2)$ entspricht einem gemeinsamen
Punkt der Geraden.
\begin{itemize}[leftmargin=1.6em,itemsep=1.2em]
  \item
    \begin{minipage}[t]{\linewidth}
    \textbf{LGS hat eine eindeutige Lösung}\par\smallskip
    \noindent
    \begin{minipage}[c]{0.64\linewidth}
      Wir betrachten das LGS
      \[
        \begin{aligned}
          \textcolor{THMBlue}{x_1+x_2}&\textcolor{THMBlue}{=2},\\
          \textcolor{black!60}{x_1-x_2}&\textcolor{black!60}{=0}.
        \end{aligned}
        \quad\Longleftrightarrow\quad
        \left[
          \begin{array}{@{}rr|r@{}}
            \textcolor{THMBlue}{1}&\textcolor{THMBlue}{1}&\textcolor{THMBlue}{2}\\
            \textcolor{black!60}{1}&\textcolor{black!60}{-1}&\textcolor{black!60}{0}
          \end{array}
        \right]
      \]
      Durch Addition beider Gleichungen folgt $2x_1=2$ und damit $x_1=1$.
      Einsetzen in eine der Gleichungen ergibt $x_2=1$.

      \smallskip
      Geometrische Betrachtung: Die Geraden besitzen genau einen
      Schnittpunkt $S=(1,1)$. Somit ist
      \[
        \mathcal{L}=\Cloze{5cm}{\{(1,1)\}}.
      \]
    \end{minipage}\hfill
    \begin{minipage}[c]{0.32\linewidth}
      \centering
      \CachedTikzFigure{content/resources/chapters/linear_systems/tikz/unique_solution_lines.tex}
    \end{minipage}
    \end{minipage}

  \item
    \begin{minipage}[t]{\linewidth}
    \textbf{LGS hat keine Lösung}\par\smallskip
    \noindent
    \begin{minipage}[c]{0.64\linewidth}
      Wir betrachten das LGS
      \[
        \begin{aligned}
          \textcolor{THMBlue}{x_1+x_2}&\textcolor{THMBlue}{=2},\\
          \textcolor{black!60}{x_1+x_2}&\textcolor{black!60}{=1}.
        \end{aligned}
        \quad\Longleftrightarrow\quad
        \left[
          \begin{array}{@{}rr|r@{}}
            \textcolor{THMBlue}{1}&\textcolor{THMBlue}{1}&\textcolor{THMBlue}{2}\\
            \textcolor{black!60}{1}&\textcolor{black!60}{1}&\textcolor{black!60}{1}
          \end{array}
        \right]
      \]
      Die Gleichungen widersprechen sich gegenseitig und können daher
      nicht gleichzeitig erfüllt sein.

      \smallskip
      Geometrische Betrachtung: Die Geraden sind parallel und besitzen
      keinen Schnittpunkt. Daher gilt
      \[
        \mathcal{L}=\Cloze{5cm}{\emptyset}.
      \]
    \end{minipage}\hfill
    \begin{minipage}[c]{0.32\linewidth}
      \centering
      \CachedTikzFigure{content/resources/chapters/linear_systems/tikz/no_solution_lines.tex}
    \end{minipage}
    \end{minipage}

  \item
    \begin{minipage}[t]{\linewidth}
    \textbf{LGS hat $\infty$-viele Lösungen}\par\smallskip
    \noindent
    \begin{minipage}[c]{0.64\linewidth}
      Wir betrachten das LGS
      \[
        \begin{aligned}
          \textcolor{THMBlue}{x_1+x_2}&\textcolor{THMBlue}{=2},\\
          \textcolor{olive!80!black}{2x_1+2x_2}&\textcolor{olive!80!black}{=4}.
        \end{aligned}
        \quad\Longleftrightarrow\quad
        \left[
          \begin{array}{@{}rr|r@{}}
            \textcolor{THMBlue}{1}&\textcolor{THMBlue}{1}&\textcolor{THMBlue}{2}\\
            \textcolor{olive!80!black}{2}&\textcolor{olive!80!black}{2}&\textcolor{olive!80!black}{4}
          \end{array}
        \right]
      \]
      Die zweite Gleichung ist das Doppelte der ersten und stellt daher
      keine zusätzliche Bedingung dar. Jedes Tupel,
      das die erste Gleichung erfüllt, erfüllt auch die zweite.

      \smallskip
      Geometrische Betrachtung: Die Geraden sind identisch und besitzen
      unendlich viele Schnittpunkte. Das bedeutet nicht, dass jedes Tupel aus
      $\R^2$ eine Lösung ist: Lösungen sind genau die Punkte auf der gemeinsamen
      Geraden. Daher gilt
      \[
        \mathcal{L}=\Cloze{7.5cm}{\{(x_1,x_2)\in\R^2\mid x_1+x_2=2\}}.
      \]
    \end{minipage}\hfill
    \begin{minipage}[c]{0.32\linewidth}
      \centering
      \CachedTikzFigure{content/resources/chapters/linear_systems/tikz/infinitely_many_solutions_lines.tex}
    \end{minipage}
    \end{minipage}
\end{itemize}

\vspace{\baselineskip}

Die geometrische Betrachtung macht anschaulich, dass lediglich diese drei Fälle
auftreten können: Zwei Geraden schneiden sich, verlaufen parallel oder liegen
aufeinander.

Die in $2\mathrm{D}$ gewonnenen Erkenntnisse lassen sich auf beliebige
$m\times n$-LGS verallgemeinern.
\begin{TheoremBox}{Anzahl der Lösungen eines LGS}{linear_systems_three_solution_types}
  Jedes $m\times n$-LGS über $\R$ besitzt genau einen der drei Lösungstypen:
  \SolutionTypeCases
  Insbesondere kann ein LGS weder genau zwei noch allgemein endlich viele
  Lösungen größer als eins besitzen.
\end{TheoremBox}

\subsection{Homogene lineare Gleichungssysteme}
\label{subsec:linear_systems_homogeneous_systems}

Ein wichtiger Spezialfall entsteht, wenn sämtliche Absolutglieder eines LGS
null sind.

\begin{DefinitionBox}{Homogenes lineares Gleichungssystem}{linear_systems_homogeneous_system}
  Ein lineares Gleichungssystem heißt \emph{homogen}
  \EnglishTerm{homogeneous}, wenn alle Absolutglieder gleich $\MathHighlight{\text{null}}$ sind:
  \[
    \begin{aligned}
      a_{11}x_1+a_{12}x_2+\dots+a_{1n}x_n&=\MathHighlight{0},\\
      a_{21}x_1+a_{22}x_2+\dots+a_{2n}x_n&=\MathHighlight{0},\\
      &\mathrel{\phantom{=}}\vdots\\
      a_{m1}x_1+a_{m2}x_2+\dots+a_{mn}x_n&=\MathHighlight{0}.
    \end{aligned}
  \]
  Seine erweiterte Koeffizientenmatrix besitzt die Form $[A\mid 0]$.
\end{DefinitionBox}

\begin{TheoremBox}{Triviale Lösung homogener LGS}{linear_systems_homogeneous_trivial_solution}
  Jedes homogene LGS besitzt die \emph{triviale Lösung}
  \EnglishTerm{trivial solution}
  \[
    (x_1,\dots,x_n)=\Cloze{4.5cm}{(0,\dots,0)}.
  \]
  Denn beim Einsetzen dieser Werte wird jede Gleichung zu $0=0$.
\end{TheoremBox}

\begin{KeyPointBox}{Lösungsmengen homogener LGS}{linear_systems_homogeneous_solution_sets}
  Da die triviale Lösung immer vorhanden ist, kann die Lösungsmenge eines
  homogenen LGS nicht leer sein:
  \SolutionTypeCases*
  Ein homogenes LGS besitzt daher entweder genau eine oder unendlich viele
  Lösungen.
\end{KeyPointBox}

\subsection{Rang und Lösungstyp}
\label{subsec:linear_systems_solution_criteria}

Mit Hilfe der Zeilenstufenform kann man ablesen, zu welchem Lösungstyp
ein LGS gehört. Eine zentrale Rolle nimmt dabei der Begriff des Rangs ein.

\begin{DefinitionBox}{Rang einer Matrix}{linear_systems_matrix_rank}
  Der \emph{Rang} \EnglishTerm{rank} einer Matrix ist die Anzahl der
  Nichtnullzeilen in einer Zeilenstufenform, wobei eine Nichtnullzeile
  mindestens einen Eintrag ungleich $0$ enthält. Man schreibt
  $\operatorname{rang}(A)$.
  \begin{center}
    \CachedTikzFigure{content/resources/chapters/linear_systems/tikz/row_echelon_rank.tex}
  \end{center}
\end{DefinitionBox}

Auf diese Weise können bei einem LGS sowohl der Rang der Koeffizientenmatrix
$A$ als auch der Rang der erweiterten Koeffizientenmatrix $[A\mid b]$
berechnet werden.

Elementare Zeilenumformungen verändern den Rang nicht. Obwohl eine Matrix auf
verschiedenen Wegen in Zeilenstufenform gebracht werden kann, ist diese Anzahl
daher eindeutig.

\begin{ExampleBox}{Rang ablesen}{linear_systems_rank_example}
  \begin{enumerate}[label=\alph*),leftmargin=1.8em,itemsep=1.3em]
    \item
      \[
        [A\mid b]\sim
        \left[
          \begin{array}{@{}rrr|r@{}}
            \phantom{-}1&\phantom{-}2&\phantom{-}0&\phantom{-}3\\
            \phantom{-}0&\phantom{-}0&\phantom{-}1&\phantom{-}5
          \end{array}
        \right]
        \quad\Longrightarrow\quad
        \begin{aligned}
          n&=\Cloze{1cm}{3},\\[0.5em]
          \operatorname{rang}(A)&=\Cloze{1cm}{2},\\[0.5em]
          \operatorname{rang}([A\mid b])&=\Cloze{1cm}{2}
        \end{aligned}
      \]

    \item
      \[
        [A\mid b]\sim
        \left[
          \begin{array}{@{}rr|r@{}}
            \phantom{-}1&-1&\phantom{-}3\\
            \phantom{-}0&\phantom{-}1&\phantom{-}5\\
            \phantom{-}0&\phantom{-}0&\phantom{-}1
          \end{array}
        \right]
        \quad\Longrightarrow\quad
        \begin{aligned}
          n&=\Cloze{1cm}{2},\\[0.5em]
          \operatorname{rang}(A)&=\Cloze{1cm}{2},\\[0.5em]
          \operatorname{rang}([A\mid b])&=\Cloze{1cm}{3}
        \end{aligned}
      \]
  \end{enumerate}
\end{ExampleBox}

\begin{TheoremBox}{Lösbarkeits- und Rangkriterium}{linear_systems_rank_criterion}
  Für ein LGS mit $n$ Unbekannten, Koeffizientenmatrix $A$ und erweiterter
  Koeffizientenmatrix $[A\mid b]$ gilt:
  \[
    \begin{array}{rcll}
      \operatorname{rang}(A)
      &<&\operatorname{rang}([A\mid b])
      &\Longleftrightarrow\quad \text{keine Lösung},\\[0.7em]
      \operatorname{rang}(A)
      &=&\operatorname{rang}([A\mid b])=n
      &\Longleftrightarrow\quad \text{genau eine Lösung},\\[0.7em]
      \operatorname{rang}(A)
      &=&\operatorname{rang}([A\mid b])<n
      &\Longleftrightarrow\quad \infty\text{-viele Lösungen}.
    \end{array}
  \]
\end{TheoremBox}

Der Rangvergleich prüft zunächst die \emph{Verträglichkeit}
\EnglishTerm{consistency} der Gleichungen. Sind die Ränge gleich, entscheidet
der Vergleich von $\operatorname{rang}(A)$ mit der Zahl $n$ der Unbekannten,
ob genau eine oder $\infty$-viele Lösungen vorliegen.

\clearpage
\begin{ExampleBox}{Lösungstyp an der Zeilenstufenform erkennen}{linear_systems_rank_solution_types_example}
  \begin{enumerate}[label=\alph*),leftmargin=1.8em,itemsep=1.3em]
    \item
      \[
        [A\mid b]\sim
        \left[
          \begin{array}{@{}rr|r@{}}
            1&2&3\\
            0&0&1\\
            0&0&0
          \end{array}
        \right]
        \quad\Longrightarrow\quad
        \left.
        \begin{aligned}
          n&=\Cloze{1cm}{2},\\[0.5em]
          \operatorname{rang}(A)&=\Cloze{1cm}{1},\\[0.5em]
          \operatorname{rang}([A\mid b])&=\Cloze{1cm}{2}
        \end{aligned}
        \quad\right\}
        \quad\Longrightarrow\quad
        \Cloze{4cm}{\text{keine Lösung}}
      \]

    \item
      \[
        [A\mid b]\sim
        \left[
          \begin{array}{@{}rr|r@{}}
            1&2&3\\
            0&1&5
          \end{array}
        \right]
        \quad\Longrightarrow\quad
        \left.
        \begin{aligned}
          n&=\Cloze{1cm}{2},\\[0.5em]
          \operatorname{rang}(A)&=\Cloze{1cm}{2},\\[0.5em]
          \operatorname{rang}([A\mid b])&=\Cloze{1cm}{2}
        \end{aligned}
        \quad\right\}
        \quad\Longrightarrow\quad
        \Cloze{4cm}{\text{genau eine Lösung}}
      \]

    \item
      \[
        [A\mid b]\sim
        \left[
          \begin{array}{@{}rrr|r@{}}
            1&2&0&3\\
            0&0&1&5
          \end{array}
        \right]
        \quad\Longrightarrow\quad
        \left.
        \begin{aligned}
          n&=\Cloze{1cm}{3},\\[0.5em]
          \operatorname{rang}(A)&=\Cloze{1cm}{2},\\[0.5em]
          \operatorname{rang}([A\mid b])&=\Cloze{1cm}{2}
        \end{aligned}
        \quad\right\}
        \quad\Longrightarrow\quad
        \Cloze{4cm}{\infty\text{-viele Lösungen}}
      \]
  \end{enumerate}
\end{ExampleBox}


\begin{KeyPointBox}{Widerspruchszeile}{linear_systems_inconsistent_row}
  Ein lineare Gleich, die niemals erfüllt werden kann, ist
  \[
    0x_1+0x_2+\dots+0x_n=c
    \quad\Longleftrightarrow\quad
    \left[\begin{array}{@{}cccc|c@{}}
      0&0&\cdots&0&c
    \end{array}\right],
    \qquad c\neq0,
  \]
  und heißt \emph{Widerspruchszeile} \EnglishTerm{inconsistent row}. 
  Für ein LGS in Zeilenstufenform gilt daher
  \[
    \exists\;\text{Widerspruchszeile}
    \quad\Longleftrightarrow\quad
    \operatorname{rang}(A)<\operatorname{rang}([A\mid b])
    \quad\Longleftrightarrow\quad
    \text{keine Lösung}.
  \]
\end{KeyPointBox}


\begin{ExampleBox}{Lösungstyp eines $3\times3$-LGS bestimmen}{linear_systems_inconsistent_three_by_three_example}
  Bringen Sie die erweiterte Koeffizientenmatrix des folgenden LGS in
  Zeilenstufenform und bestimmen Sie den Lösungstyp:
  \[
    \begin{aligned}
      x_1+x_2+x_3&=3,\\
      2x_1+x_2-x_3&=1,\\
      3x_1+2x_2&=5
    \end{aligned}
    \quad\Longleftrightarrow\quad
    [A\mid b]=
    \left[
      \begin{array}{@{}rrr|r@{}}
        1&1& 1&3\\
        2&1&-1&1\\
        3&2& 0&5
      \end{array}
    \right].
  \]

  \begin{ClozeArea}[32]{35}
    \textbf{Vorwärtselimination:}
    \[
      \begin{array}{r@{\;}r@{\qquad}l}
        &\left[\begin{array}{@{}rrr|r@{}}
          1&1&1&3\\
          2&1&-1&1\\
          3&2&0&5
        \end{array}\right]
        &\begin{array}{@{}l@{}}
          {}\\
          \smash{\TransformOp{Z_2-2Z_1}}\\
          \smash{\TransformOp{Z_3-3Z_1}}
        \end{array}\\[1.0em]
        \Longleftrightarrow
        &\left[\begin{array}{@{}rrr|r@{}}
          1&1&1&3\\
          0&-1&-3&-5\\
          0&-1&-3&-4
        \end{array}\right]
        &\begin{array}{@{}l@{}}
          {}\\
          {}\\
          \smash{\TransformOp{Z_3-Z_2}}
        \end{array}\\[1.0em]
        \Longleftrightarrow
        &\left[\begin{array}{@{}rrr|r@{}}
          1&1&1&3\\
          0&-1&-3&-5\\
          \textcolor{red!70!black}{0}&\textcolor{red!70!black}{0}&
          \textcolor{red!70!black}{0}&\textcolor{red!70!black}{1}
        \end{array}\right]
        &\text{Zeilenstufenform}
      \end{array}
    \]
    Die letzte Zeile ist die
    \textcolor{red!70!black}{\emph{Widerspruchszeile}}:
    \[
      \textcolor{red!70!black}{
        0x_1+0x_2+0x_3=1
        \quad\Longleftrightarrow\quad
        0=1.
      }
    \]
    \textbf{Rangvergleich:}
    \[
      \left.
      \begin{aligned}
        n&=3,\\[0.5em]
        \operatorname{rang}(A)&=2,\\[0.5em]
        \operatorname{rang}([A\mid b])&=3
      \end{aligned}
      \quad\right\}
      \quad\Longrightarrow\quad
      \text{keine Lösung}.
    \]
    Damit ist die Lösungsmenge
    \[
      \mathcal L=\varnothing.
    \]
  \end{ClozeArea}
\end{ExampleBox}

\clearpage
\subsection{Unendlich viele Lösungen parametrisieren}
\label{subsec:linear_systems_parametric_solution_sets}

Besitzt ein lösbares LGS $\infty$-viele Lösungen, wird seine Lösungsmenge
mithilfe von Parametern beschrieben.
Das folgende Beispiel zeigt dies anhand eines bereits bekannten LGS.

\begin{ExampleBox}{Eine Lösungsgerade parametrisieren}{linear_systems_parametric_solution}
  Das folgende $2\times2$-LGS haben wir zuvor bereits geometrisch betrachtet.
  Nun untersuchen wir es zusätzlich mithilfe des Rangs und lösen es mit dem
  Gauß-Verfahren:
  \[
    \begin{aligned}
      \textcolor{THMBlue}{x_1+x_2}&\textcolor{THMBlue}{=2},\\
      \textcolor{olive!80!black}{2x_1+2x_2}&\textcolor{olive!80!black}{=4}.
    \end{aligned}
    \quad\Longleftrightarrow\quad
    \left[
      \begin{array}{@{}rr|r@{}}
        \Cloze{0.7cm}{\textcolor{THMBlue}{1}}&
        \Cloze{0.7cm}{\textcolor{THMBlue}{1}}&
        \Cloze{0.7cm}{\textcolor{THMBlue}{2}}\\[0.35em]
        \Cloze{0.7cm}{\textcolor{olive!80!black}{2}}&
        \Cloze{0.7cm}{\textcolor{olive!80!black}{2}}&
        \Cloze{0.7cm}{\textcolor{olive!80!black}{4}}
      \end{array}
    \right].
  \]
  Die Vorwärtselimination liefert die Zeilenstufenform
  \begin{ClozeArea}[32]{8}
    \[
      \begin{array}{r@{\;}r@{\qquad}l}
        &\left[\begin{array}{@{}rr|r@{}}
          \textcolor{THMBlue}{1}&\textcolor{THMBlue}{1}&\textcolor{THMBlue}{2}\\
          \textcolor{olive!80!black}{2}&\textcolor{olive!80!black}{2}&
          \textcolor{olive!80!black}{4}
        \end{array}\right]
        &\begin{array}{@{}l@{}}
          {}\\
          \smash{\TransformOp{Z_2-2Z_1}}
        \end{array}\\[1.0em]
        \Longleftrightarrow
        &\left[\begin{array}{@{}rr|r@{}}
          \textcolor{THMBlue}{1}&\textcolor{THMBlue}{1}&\textcolor{THMBlue}{2}\\
          \textcolor{olive!80!black}{0}&\textcolor{olive!80!black}{0}&
          \textcolor{olive!80!black}{0}
        \end{array}\right]
        &\text{Zeilenstufenform}
      \end{array}
    \]
  \end{ClozeArea}
  Es gilt
  \[
    \left.
    \begin{aligned}
      n&=\Cloze{1cm}{2},\\[0.5em]
      \operatorname{rang}(A)&=\Cloze{1cm}{1},\\[0.5em]
      \operatorname{rang}([A\mid b])&=\Cloze{1cm}{1}
    \end{aligned}
    \quad\right\}
    \quad\Longrightarrow\quad
    \Cloze{4cm}{\infty\text{-viele Lösungen}}.
  \]

  \noindent
  \begin{minipage}[c]{0.62\linewidth}
    Um die Lösungsmenge geschickt anzugeben, führen wir einen Parameter
    $t\in\R$ ein und setzen $x_2=t$ in Zeile 1 ein.
    \[
      \begin{aligned}
        \Cloze{10cm}{
        x_1+x_2=2
        \quad\overset{x_2\,=\,t}{\Longleftrightarrow}\quad x_1+t=2
        \quad\Longleftrightarrow\quad x_1=2-t.
        }
      \end{aligned}
    \]
    Damit lautet die Lösungsmenge
    \[
      \mathcal{L}
      =\Cloze{7.8cm}{
        \left\{
          \begin{pmatrix}2\\0\end{pmatrix}
          +t\begin{pmatrix}-1\\1\end{pmatrix}
          \;\middle|\;t\in\R
        \right\}
      }.
    \]
  \end{minipage}\hfill
  \begin{minipage}[c]{0.34\linewidth}
    \centering
    \CachedTikzFigure{content/resources/chapters/linear_systems/tikz/infinitely_many_solutions_lines.tex}
  \end{minipage}

  \smallskip
  Der Parameter $t$ wirkt wie ein Schieberegler: Für jeden Wert von $t$
  entsteht eine andere Lösung $(2-t,t)$. Durchläuft $t$ alle reellen Zahlen,
  durchläuft der zugehörige Punkt die gesamte Lösungsgerade $x_1+x_2=2$.
\end{ExampleBox}

\begin{KeyPointBox}{Anzahl der freien Variablen}{linear_systems_number_free_variables}
  Variablen für man einen freien Parameter einführt, nennt man \emph{freie Variablen} 
  \EnglishTerm{free variables}. Für ein lösbares LGS mit $n$ Unbekannten ist die Anzahl
  freier Variablen gegeben durch
  \[
    \#\{\text{freie Variablen}\}
    = \#\{\text{benötigte Parameter}\}
    =n-\operatorname{rang}(A).
  \]
  \par\smallskip
  \hfill
  {\footnotesize\color{black!55}
    $\#\{\ldots\}\mathrel{\widehat{=}}\text{Anzahl der Elemente von }\{\ldots\}$.}
\end{KeyPointBox}


Die Anzahl der freien Variablen bestimmt zugleich die Dimension der
Lösungsmenge:
\[
  \begin{aligned}
    0\ \text{freie Variable}
    &\quad\Longleftrightarrow\quad
    \dim(\mathcal L)=0 \quad \text{(eindeutiger Lösungspunkt)} \\[0.5em]
    1\ \text{freie Variable}
    &\quad\Longleftrightarrow\quad
    \dim(\mathcal L)=1 \quad \text{(Lösungsgerade)} \\[0.5em]
    2\ \text{freie Variablen}
    &\quad\Longleftrightarrow\quad
    \dim(\mathcal L)=2 \quad \text{(Lösungsebene)} \\[0.5em]
    3\ \text{freie Variablen}
    &\quad\Longleftrightarrow\quad
    \dim(\mathcal L)=3 \quad \text{(Lösungsvolumen)} \\[0.5em]
    &\hspace{-2.8em}\vdots &&
  \end{aligned}
\]


\clearpage
\section{Determinante für LGS}
\label{sec:linear_systems_determinants}

Die Determinante ist eine zentrale Größe der Mathematik. Sie spielt insbesondere
bei linearen Gleichungssystemen, in der Matrix- und Vektorrechnung sowie in der
mehrdimensionalen Analysis eine wichtige Rolle.

In diesem Abschnitt führen wir die Determinante zur Lösung linearer
Gleichungssysteme ein und lernen ihre Aussagekraft hinsichtlich des Lösungstyps
kennen.

\subsection{Determinante von $2\times2$-Matrizen}
\label{subsec:linear_systems_two_by_two_determinants}

Wir lösen nun das allgemeine $2\times2$-LGS mit dem Gauß-Verfahren, um die
Struktur seiner Lösung genauer zu untersuchen.

\begin{ExampleBox}{Allgemeines $2\times2$-LGS}{linear_systems_determinant_from_gaussian_elimination}
  Wir lösen das folgende $2\times2$-LGS mit dem Gauß-Verfahren
  \[
    \begin{aligned}
      a_{11}x_1+a_{12}x_2&=b_1,\\
      a_{21}x_1+a_{22}x_2&=b_2
    \end{aligned}
    \qquad\Longleftrightarrow\qquad
    \left[
      \begin{array}{@{}cc|c@{}}
        \Cloze{0.9cm}{a_{11}}&\Cloze{0.9cm}{a_{12}}&\Cloze{0.9cm}{b_1}\\[0.3em]
        \Cloze{0.9cm}{a_{21}}&\Cloze{0.9cm}{a_{22}}&\Cloze{0.9cm}{b_2}
      \end{array}
    \right]
  \]

  \begin{ClozeArea}[34]{24}
    \textbf{Vorwärtselimination:}
    \[
      \begin{array}{r@{\;}r@{\qquad}l}
        &\left[
          \begin{array}{@{}cc|c@{}}
            a_{11}&a_{12}&b_1\\
            a_{21}&a_{22}&b_2
          \end{array}
        \right]
        &\raisebox{-0.65em}{\(\TransformOp{a_{11}Z_2-a_{21}Z_1}\)}\\[1.2em]
        \Longleftrightarrow
        &\left[
          \begin{array}{@{}cc|c@{}}
            a_{11}&a_{12}&b_1\\
            0&a_{11}a_{22}-a_{12}a_{21}&a_{11}b_2-a_{21}b_1
          \end{array}
        \right]
        &\text{Zeilenstufenform}
      \end{array}
    \]
  \end{ClozeArea}
\end{ExampleBox}

\clearpage
\begin{tcolorbox}[
  scriptbox,
  breakable=false,
  borderline west={1.2pt}{0pt}{black!40},
  title={Beispiel~\ref{ex:linear_systems_determinant_from_gaussian_elimination}: Allgemeines $2\times2$-LGS -- Fortsetzung}
]
  \begin{ClozeArea}[34]{22}
    \textbf{Rückwärtseinsetzen:} Für
    $a_{11}a_{22}-a_{12}a_{21}\neq0$ folgt
    \begin{itemize}[leftmargin=2.1em,itemsep=0.7em]
      \item In Zeile 2 einsetzen:
        \[
          (a_{11}a_{22}-a_{12}a_{21})x_2
          =a_{11}b_2-a_{21}b_1
          \quad\Longleftrightarrow\quad
          x_2=\frac{a_{11}b_2-a_{21}b_1}
          {a_{11}a_{22}-a_{12}a_{21}}.
        \]
      \item in Zeile 1 einsetzen:
        \[
          \begin{aligned}
            a_{11}x_1+a_{12}x_2=b_1
            \quad\Longleftrightarrow\quad&
            a_{11} x_1 + a_{12} \frac{a_{11}b_2-a_{21}b_1}{a_{11}a_{22}-a_{12}a_{21}} = b_1 ,\\[0.6em]
            \quad\Longleftrightarrow\quad&
            x_1=\frac{a_{22}b_1-a_{12}b_2}{a_{11}a_{22}-a_{12}a_{21}}.
          \end{aligned}
        \]
    \end{itemize}
    Damit ergibt sich die folgende Lösungsmenge:
    \[
      \mathcal L=
      \left\{
        \left(
          \frac{a_{22}b_1-a_{12}b_2}{a_{11}a_{22}-a_{12}a_{21}},
          \frac{a_{11}b_2-a_{21}b_1}{a_{11}a_{22}-a_{12}a_{21}}
        \right)
      \right\}.
    \]
  \end{ClozeArea}
\end{tcolorbox}

Die Ergebnisse des LGS lassen sich mit der folgenden Definition einer Determinante
bedeutend leichter und übersichtlicher schreiben.

\begin{DefinitionBox}{Determinante einer $2\times2$-Matrix}{linear_systems_two_by_two_determinant}
  Für eine $2\times2$ Koeffizientenmatrix
  \[
    A=\begin{pmatrix}
      \textcolor{THMBlue}{a}&\textcolor{olive!80!black}{b}\\
      \textcolor{olive!80!black}{c}&\textcolor{THMBlue}{d}
    \end{pmatrix}
  \]
  berechnet sich die \emph{Determinante} \EnglishTerm{determinant} von $A$ durch
  \[
    \det(A)
    =\det\!\begin{pmatrix}
      \textcolor{THMBlue}{a}&\textcolor{olive!80!black}{b}\\
      \textcolor{olive!80!black}{c}&\textcolor{THMBlue}{d}
    \end{pmatrix}
    =\begin{vmatrix}
      \textcolor{THMBlue}{a}&\textcolor{olive!80!black}{b}\\
      \textcolor{olive!80!black}{c}&\textcolor{THMBlue}{d}
    \end{vmatrix}
    =\textcolor{THMBlue}{ad}-\textcolor{olive!80!black}{bc}
  \]
  \textbf{Merkhilfe:}
  \enquote{\textcolor{THMBlue}{Hauptdiagonale} minus
  \textcolor{olive!80!black}{Gegendiagonale}.}
\end{DefinitionBox}

\begin{ExampleBox}{Determinante einer $2\times2$-Matrix}{linear_systems_two_by_two_determinant_calculation}
  \begin{enumerate}[label=\alph*),leftmargin=1.8em,itemsep=1.0em]
    \item
      \(
        \det\!\begin{pmatrix}3&2\\1&4\end{pmatrix}
        =\Cloze{6.5cm}{3\cdot4-2\cdot1=10}.
      \)
    \item
      \(
        \det\!\begin{pmatrix}2&-1\\5&3\end{pmatrix}
        =\Cloze{6.5cm}{2\cdot3-(-1)\cdot5=11}.
      \)
  \end{enumerate}
\end{ExampleBox}

\begin{samepage}
Mit der Definition der $2\times2$ Determinante lassen sich die Lösungen des allgemeinen
$2\times2$ LGS aus Beispiel~\ref{ex:linear_systems_determinant_from_gaussian_elimination}
nun sehr übersichtlich schreiben.
\[
  x_1=
  \frac{
    \det\!\begin{pmatrix}
      \textcolor{THMBlue}{b_1}&a_{12}\\
      \textcolor{THMBlue}{b_2}&a_{22}
    \end{pmatrix}
  }{
    \det\!\begin{pmatrix}a_{11}&a_{12}\\a_{21}&a_{22}\end{pmatrix}
  },
  \qquad
  x_2=
  \frac{
    \det\!\begin{pmatrix}
      a_{11}&\textcolor{THMBlue}{b_1}\\
      a_{21}&\textcolor{THMBlue}{b_2}
    \end{pmatrix}
  }{
    \det\!\begin{pmatrix}a_{11}&a_{12}\\a_{21}&a_{22}\end{pmatrix}
  }.
\]
\end{samepage}

\subsection{Die Cramersche Regel}
\label{subsec:linear_systems_cramers_rule}

Die Formeln für $x_1$ und $x_2$ besitzen ein auffällige Struktur:
\begin{itemize}
  \item Im Nenner steht jeweils die Determinante der Koeffizientenmatrix.
  \item Im Zähler steht die Determinante einer Matrix, in der die zur
  gesuchten Unbekannten gehörende Spalte durch die Absolutglieder ersetzt wurde.
\end{itemize}
Für eine kompakte Schreibweise führen wir dafür die Cramerschen Ersatzmatrizen ein.

\begin{DefinitionBox}{Cramersche Ersatzmatrix}{linear_systems_cramer_replacement_matrix}
  Gegeben sei ein $n\times n$-LGS mit Koeffizientenmatrix $A$ und dem Vektor
  der Absolutglieder $b$. Ersetzt man in $A$ die $i$-te Spalte durch $b$, so
  entsteht die $i$-te \emph{Cramersche Ersatzmatrix}
  \EnglishTerm{Cramer replacement matrix} $A_i$.
  \par\medskip
  Beispielhaft für den $2\times2$-Fall:
  \begin{center}
    \CachedTikzFigure{content/resources/chapters/linear_systems/tikz/cramers_column_replacement.tex}
  \end{center}
\end{DefinitionBox}

Mit den Erstzmatrizen $A_i$ lässt sich die Lösung des allgemeinen $2 \times 2$ LGS
nochmals kompakter und übersichtlicher formulieren.
%
\begin{samepage}
Für ein $2\times2$-LGS mit $\det(A)\neq0$ lauten die Lösungen damit
\[
  x_1=\frac{\det(A_1)}{\det(A)},
  \qquad
  x_2=\frac{\det(A_2)}{\det(A)}.
\]

Die Cramersche Regel verallgemeinert dieses Prinzip auf beliebig viele Unbekannte.
\begin{TheoremBox}{Cramersche Regel}{linear_systems_cramers_rule_two_by_two}
  Für ein $n\times n$-LGS mit Koeffizientenmatrix $A$ und $\det(A)\neq0$
  gilt
  \[
    x_i=\frac{\det(A_i)}{\det(A)},
    \qquad i=1,\dots,n.
  \]
  Diese Lösungsformel heißt \emph{Cramersche Regel}
  \EnglishTerm{Cramer's rule}.
\end{TheoremBox}
\end{samepage}

\begin{ExampleBox}{Gauß-Verfahren versus Cramersche Regel}{linear_systems_cramers_rule_example}
  Wir lösen das folgende LGS mit 2 verschiedenen Verfahren.
  \[
    \begin{aligned}
      x_1+x_2&=2,\\
      x_1-x_2&=-4.
    \end{aligned}
    \qquad\Longleftrightarrow\qquad
    [A\mid b]=
    \left[
      \begin{array}{@{}rr|r@{}}
        1& 1& 2\\
        1&-1&-4
      \end{array}
    \right]
  \]
  {\setlength{\tabcolsep}{0.1cm}%
  \setlength{\arrayrulewidth}{0.7pt}%
  \noindent\begin{tabularx}{\linewidth}{@{}X|X@{}}
  \centering\textbf{Gauß-Verfahren}
  &
  \centering\arraybackslash\textbf{Cramersche Regel}\\[0.2em]
  \hline
  \begin{minipage}[t]{\linewidth}
    \vspace{0.45em}
    \begin{ClozeArea}[16]{22}
      Vorwärtselimination:
      \[
        \begin{array}{r@{\;}c@{\qquad}l}
          &\left[\begin{array}{@{}rr|r@{}}
            1& 1& 2\\
            1&-1&-4
          \end{array}\right]
          &\raisebox{-0.65em}{\(\TransformOp{Z_2-Z_1}\)}\\[1.0em]
          \Longleftrightarrow
          &\left[\begin{array}{@{}rr|r@{}}
            1& 1& 2\\
            0&-2&-6
          \end{array}\right]
          &\text{Zeilenstufenform}
        \end{array}
      \]
      \smallskip
      Rückwärtseinsetzen:
      \begin{itemize}[leftmargin=1.4em,itemsep=0.3em,topsep=0.2em]
        \item In Zeile 2 einsetzen:
          \[
            -2x_2=-6
            \quad\Longleftrightarrow\quad x_2=3.
          \]
        \item In Zeile 1 einsetzen:
          \[
            \begin{aligned}
              x_1+x_2=2
              &\quad\Longleftrightarrow\quad x_1+3=2,\\
              &\quad\Longleftrightarrow\quad x_1=-1.
            \end{aligned}
          \]
      \end{itemize}
      Lösungsmenge:
      \[
        \mathcal L=\{(-1,3)\}.
      \]
    \end{ClozeArea}
  \end{minipage}
  &
  \begin{minipage}[t]{\linewidth}
    \vspace{0.45em}
    \begin{ClozeArea}[16]{22}
      \[
        A=\begin{pmatrix}1&1\\1&-1\end{pmatrix}.
      \]
      Cramersche Ersatzmatrizen:
      \[
        A_1=\begin{pmatrix}2&1\\-4&-1\end{pmatrix},
        \qquad
        A_2=\begin{pmatrix}1&2\\1&-4\end{pmatrix}.
      \]
      Determinanten:
      \[
        \begin{aligned}
          \det(A)
            &=\det\!\begin{pmatrix}1&1\\1&-1\end{pmatrix}
              =1\cdot(-1)-1\cdot1=-2,\\[0.8em]
          \det(A_1)
            &=\det\!\begin{pmatrix}2&1\\-4&-1\end{pmatrix}
              =2\cdot(-1)-1\cdot(-4)=2,\\[0.8em]
          \det(A_2)
            &=\det\!\begin{pmatrix}1&2\\1&-4\end{pmatrix}
              =1\cdot(-4)-2\cdot1=-6.
        \end{aligned}
      \]
      Cramersche Regel:
      \[
        x_1=\frac{\det(A_1)}{\det(A)}
        =\frac{2}{-2}=-1,
      \]
      \[
        x_2=\frac{\det(A_2)}{\det(A)}
        =\frac{-6}{-2}=3.
      \]
      Lösungsmenge:
      \[
        \mathcal L=\{(-1,3)\}.
      \]
    \end{ClozeArea}
  \end{minipage}
  \end{tabularx}}
  \par\smallskip
  Mit beiden Verfahren wird dieselbe Lösungsmenge ermittelt.
\end{ExampleBox}

Um die Cramersche Regel auch für allgemeine $n \times n$ lineare Gleichungssysteme
anwenden zu können, müssen wir zunächst lernen, wie eine $n \times n$-Determinante
berechnet wird.

\subsection{Determinante von $n\times n$-Matrizen}
\label{subsec:linear_systems_laplace_expansion}

Eine Determinante höherer Ordnung lässt sich schrittweise auf kleinere,
bereits berechenbare Determinanten zurückführen. Dieses Verfahren heißt
\emph{Laplace-Entwicklung} \EnglishTerm{Laplace expansion}.

\begin{TheoremBox}{Kochrezept: Laplace-Entwicklung}{linear_systems_laplace_expansion}
\begin{enumerate}[label=\textbf{Schritt \arabic*:},leftmargin=*,itemsep=0.15em]
  \item \textbf{Gedachte Vorzeichen.}
    Die gedachten Vorzeichen werden nach dem Schachbrettmuster ergänzt.
    Sie gehören nicht zur Matrix:
    \begin{center}
      \vspace{-0.35\baselineskip}
      \CachedTikzFigure{content/resources/chapters/linear_systems/tikz/laplace_sign_pattern_start.tex}
      \vspace{-0.35\baselineskip}
    \end{center}
  \item \textbf{Zeile oder Spalte auswählen.}
    Für die Entwicklung wird eine beliebige Zeile oder Spalte ausgewählt;
    das Ergebnis ist unabhängig von dieser Wahl:

    \begin{center}
      \vspace{-0.2\baselineskip}
      \CachedTikzFigure{content/resources/chapters/linear_systems/tikz/laplace_row_or_column_selection.tex}
      \vspace{-0.2\baselineskip}
    \end{center}
  \item \textbf{Entwicklungsschritt:}
    Wir entwicklen die Determinate in die folgende Summe entlang der gewählten Zeile oder Spalte:
    \[
      \det(A)
      =\sum_{\substack{\text{Einträge der gewählten}\\
                       \text{Zeile oder Spalte}}}
       \textcolor{olive!80!black}{\text{„gedachtes Vorzeichen“}}
       \;\cdot\;\textcolor{THMBlue}{a_{ij}}
       \;\cdot\;\det(A_{ij}).
    \]
    Dabei ist $A_{ij}$ die sogenannte \emph{Untermatrix} \EnglishTerm{submatrix}, die
    sich aus Matrix $A$ durch Streichen von Zeile $i$ und Spalte $j$ ergibt.

    \begin{center}
      \vspace{-0.25\baselineskip}
      \CachedTikzFigure{content/resources/chapters/linear_systems/tikz/laplace_submatrix_construction.tex}
      \vspace{-0.25\baselineskip}
    \end{center}

\end{enumerate}
Durch die Laplace-Entwicklung wird eine $n\times n$-Determinante auf $n$
Determinanten der Größe $(n-1)\times(n-1)$ zurückgeführt.
\end{TheoremBox}

\begin{minipage}{\linewidth}
Mithilfe der Laplace-Entwicklung führen wir $n\times n$-Determinanten
schrittweise auf Determinanten niedrigerer Ordnung zurück, bis nur noch
$2\times2$-Determinanten verbleiben. Diese können wir nach
Definition~\ref{def:linear_systems_two_by_two_determinant} berechnen.

\begin{center}
  \CachedTikzFigure{content/resources/chapters/linear_systems/tikz/laplace_expansion_flow.tex}
\end{center}
\end{minipage}

\begin{ExampleBox}{Laplace-Entwicklung einer $3\times3$-Determinante}{linear_systems_three_by_three_determinant}
  Wir berechnen die Determinante mit der Laplace-Entwicklung einmal nach der
  ersten Zeile und einmal nach der zweiten Spalte.
  \[
    D=\det\!\begin{pmatrix}
      {}^{\textcolor{olive!80!black}{\boldsymbol{+}}}\!1&
      {}^{\textcolor{olive!80!black}{\boldsymbol{-}}}\!2&
      {}^{\textcolor{olive!80!black}{\boldsymbol{+}}}\!3\\[0.35em]
      {}^{\textcolor{olive!80!black}{\boldsymbol{-}}}\!4&
      {}^{\textcolor{olive!80!black}{\boldsymbol{+}}}\!5&
      {}^{\textcolor{olive!80!black}{\boldsymbol{-}}}\!6\\[0.35em]
      {}^{\textcolor{olive!80!black}{\boldsymbol{+}}}\!7&
      {}^{\textcolor{olive!80!black}{\boldsymbol{-}}}\!8&
      {}^{\textcolor{olive!80!black}{\boldsymbol{+}}}\!0
    \end{pmatrix}.
  \]
  \begin{enumerate}[label=\alph*),leftmargin=1.8em,itemsep=0.8em]
    \item \textbf{Entwicklung nach der ersten Zeile:}
      \begin{ClozeArea}[32]{9}
        \[
          \begin{aligned}
            D
            &=\det\!\begin{pNiceMatrix}
                \CodeBefore
                  \rowcolor{yellow!55}{1}
                \Body
                {}^{\textcolor{olive!80!black}{\pmb{+}}}\!1&
                {}^{\textcolor{olive!80!black}{\pmb{-}}}\!2&
                {}^{\textcolor{olive!80!black}{\pmb{+}}}\!3\\[0.35em]
                {}^{\textcolor{olive!80!black}{\pmb{-}}}\!4&
                {}^{\textcolor{olive!80!black}{\pmb{+}}}\!5&
                {}^{\textcolor{olive!80!black}{\pmb{-}}}\!6\\[0.35em]
                {}^{\textcolor{olive!80!black}{\pmb{+}}}\!7&
                {}^{\textcolor{olive!80!black}{\pmb{-}}}\!8&
                {}^{\textcolor{olive!80!black}{\pmb{+}}}\!0
              \end{pNiceMatrix}
              \;=\;\textcolor{olive!80!black}{\boldsymbol{+}}\,1\det\!\begin{pmatrix}5&6\\8&0\end{pmatrix}
              \textcolor{olive!80!black}{\boldsymbol{-}}\,2\det\!\begin{pmatrix}4&6\\7&0\end{pmatrix}
              \textcolor{olive!80!black}{\boldsymbol{+}}\,3\det\!\begin{pmatrix}4&5\\7&8\end{pmatrix}\\[0.7em]
            &=1(5\cdot0-6\cdot8)
              -2(4\cdot0-6\cdot7)
              +3(4\cdot8-5\cdot7)\\[0.7em]
            &=-48+84-9=27.
          \end{aligned}
        \]
      \end{ClozeArea}
    \item \textbf{Entwicklung nach der zweiten Spalte:}
      \begin{ClozeArea}[32]{9}
        \[
          \begin{aligned}
            D
            &=\det\!\begin{pNiceMatrix}
                \CodeBefore
                  \columncolor{yellow!55}{2}
                \Body
                {}^{\textcolor{olive!80!black}{\pmb{+}}}\!1&
                {}^{\textcolor{olive!80!black}{\pmb{-}}}\!2&
                {}^{\textcolor{olive!80!black}{\pmb{+}}}\!3\\[0.35em]
                {}^{\textcolor{olive!80!black}{\pmb{-}}}\!4&
                {}^{\textcolor{olive!80!black}{\pmb{+}}}\!5&
                {}^{\textcolor{olive!80!black}{\pmb{-}}}\!6\\[0.35em]
                {}^{\textcolor{olive!80!black}{\pmb{+}}}\!7&
                {}^{\textcolor{olive!80!black}{\pmb{-}}}\!8&
                {}^{\textcolor{olive!80!black}{\pmb{+}}}\!0
              \end{pNiceMatrix}
              \;=\;\textcolor{olive!80!black}{\boldsymbol{-}}\,2\det\!\begin{pmatrix}4&6\\7&0\end{pmatrix}
              \textcolor{olive!80!black}{\boldsymbol{+}}\,5\det\!\begin{pmatrix}1&3\\7&0\end{pmatrix}
              \textcolor{olive!80!black}{\boldsymbol{-}}\,8\det\!\begin{pmatrix}1&3\\4&6\end{pmatrix}\\[0.7em]
            &=-2(4\cdot0-6\cdot7)
              +5(1\cdot0-3\cdot7)
              -8(1\cdot6-3\cdot4)\\[0.7em]
            &=84-105+48=27.
          \end{aligned}
        \]
      \end{ClozeArea}
  \end{enumerate}
\end{ExampleBox}

\begin{minipage}{\linewidth}
Bei der Laplace Entwicklung können sich durch die geschickte Wahl der Entwicklungszeile
oder -spalte auch aufwendige Rechnungen drastisch vereinfachen.

\begin{ExampleBox}{Laplace-Entwicklung einer dünn besetzten Matrix}{linear_systems_laplace_sparse_column}
  Eine \emph{dünn besetzte Matrix} \EnglishTerm{sparse matrix} enthält
  viele Nulleinträge. Wir reduzieren den Rechenaufwand deutlich, wenn wir nach Zeilen
  oder Spalten mit vielen Nulleinträgen entwickeln.
  \begin{ClozeArea}[32]{11}
        \[
          \begin{aligned}
            D
            &=\det\!\begin{pNiceMatrix}
                \CodeBefore
                  \columncolor{yellow!55}{3}
                \Body
                {}^{\textcolor{olive!80!black}{\pmb{+}}}\!1&
                {}^{\textcolor{olive!80!black}{\pmb{-}}}\!2&
                {}^{\textcolor{olive!80!black}{\pmb{+}}}\!0&
                {}^{\textcolor{olive!80!black}{\pmb{-}}}\!4\\[0.3em]
                {}^{\textcolor{olive!80!black}{\pmb{-}}}\!0&
                {}^{\textcolor{olive!80!black}{\pmb{+}}}\!7&
                {}^{\textcolor{olive!80!black}{\pmb{-}}}\!0&
                {}^{\textcolor{olive!80!black}{\pmb{+}}}\!2\\[0.3em]
                {}^{\textcolor{olive!80!black}{\pmb{+}}}\!0&
                {}^{\textcolor{olive!80!black}{\pmb{-}}}\!5&
                {}^{\textcolor{olive!80!black}{\pmb{+}}}\!0&
                {}^{\textcolor{olive!80!black}{\pmb{-}}}\!4\\[0.3em]
                {}^{\textcolor{olive!80!black}{\pmb{-}}}\!3&
                {}^{\textcolor{olive!80!black}{\pmb{+}}}\!3&
                {}^{\textcolor{olive!80!black}{\pmb{-}}}\!(-1)&
                {}^{\textcolor{olive!80!black}{\pmb{+}}}\!2
              \end{pNiceMatrix}\\[0.7em]
            &=\textcolor{olive!80!black}{\pmb{-}}\,(-1) \cdot
              \det\!\begin{pNiceMatrix}
                \CodeBefore
                  \columncolor{yellow!55}{1}
                \Body
                {}^{\textcolor{olive!80!black}{\pmb{+}}}\!1&
                {}^{\textcolor{olive!80!black}{\pmb{-}}}\!2&
                {}^{\textcolor{olive!80!black}{\pmb{+}}}\!4\\[0.3em]
                {}^{\textcolor{olive!80!black}{\pmb{-}}}\!0&
                {}^{\textcolor{olive!80!black}{\pmb{+}}}\!7&
                {}^{\textcolor{olive!80!black}{\pmb{-}}}\!2\\[0.3em]
                {}^{\textcolor{olive!80!black}{\pmb{+}}}\!0&
                {}^{\textcolor{olive!80!black}{\pmb{-}}}\!5&
                {}^{\textcolor{olive!80!black}{\pmb{+}}}\!4
              \end{pNiceMatrix}
              \quad\textcolor{olive!80!black}{%
                \longleftarrow\ \text{neues Schachbrettmuster:}%
              }\\[0.7em]
            &=1\det\!\begin{pmatrix}7&2\\5&4\end{pmatrix}
              =7\cdot4-2\cdot5=18.
          \end{aligned}
        \]
  \end{ClozeArea}
\end{ExampleBox}
\end{minipage}

\begin{ExampleBox}{Rechenaufwand einer $10\times10$-Determinante}{linear_systems_laplace_computational_effort}
  In wie viele $2\times2$-Determinanten zerfällt eine
  $10\times10$-Determinante bei vollständiger Laplace-Entwicklung, wenn keine
  vereinfachenden Nullen genutzt werden können?
  \begin{ClozeArea}[32]{8}
    \[
      \begin{aligned}
        1\text{ Determinante der Größe }10\times10
        &\;\longrightarrow\;10\text{ Determinanten der Größe }9\times9,\\[0.35em]
        &\;\longrightarrow\;10\cdot9\text{ Determinanten der Größe }8\times8,\\[0.35em]
        &\;\longrightarrow\;10\cdot9\cdot8\text{ Determinanten der Größe }7\times7,\\[-0.1em]
        &\hspace{3.8em}\vdots\\[-0.1em]
        &\;\longrightarrow\;10\cdot9\cdots3
          \text{ Determinanten der Größe }2\times2.
      \end{aligned}
    \]
    Somit müssen
    \[
      10\cdot9\cdots3=\frac{10!}{2}=1\,814\,400
    \]
    Determinanten der Größe $2\times2$ berechnet werden.
  \end{ClozeArea}
\end{ExampleBox}

\begin{KeyPointBox}{Rechenaufwand der Laplace-Entwicklung}{linear_systems_laplace_computational_effort}
  Bei der vollständigen Laplace-Entwicklung einer $n\times n$-Determinante
  entstehen ohne Vereinfachungen
  \[
    n\cdot(n-1)\cdots3=\frac{n!}{2}
  \]
  Determinanten der Größe $2\times2$. Der Rechenaufwand wächst somit
  faktoriell und wird für große Matrizen sehr schnell unpraktikabel.
\end{KeyPointBox}

\begin{minipage}{\linewidth}
\subsubsection{Spezialfall: Regel von Sarrus}
\label{subsubsec:linear_systems_sarrus_rule}

Für $3\times3$-Determinanten existiert neben der Laplace-Entwicklung ein
spezielles Rechenschema. Dieses ist häufig schneller anzuwenden, lässt sich
jedoch nicht auf Determinanten höherer Ordnung übertragen.

\begin{KeyPointBox}{Regel von Sarrus}{linear_systems_sarrus_rule}
  Im Spezialfall einer $3\times3$-Matrix lässt sich die Determinante auch mit
  der \emph{Regel von Sarrus} \EnglishTerm{Sarrus' rule} berechnen:

  \begin{center}
    \CachedTikzFigure{content/resources/chapters/linear_systems/tikz/sarrus_rule.tex}
  \end{center}
  \vspace{-0.3\baselineskip}
  \begin{align*}
    \det(A)
    ={}&\textcolor{THMBlue}{a_{11}a_{22}a_{33}
       +a_{12}a_{23}a_{31}
       +a_{13}a_{21}a_{32}}\\
     &\textcolor{red!80!black}{-a_{13}a_{22}a_{31}
       -a_{11}a_{23}a_{32}
       -a_{12}a_{21}a_{33}}.
  \end{align*}
  \textbf{Achtung:} Die Regel von Sarrus gilt ausschließlich für
  $3\times3$-Matrizen.
\end{KeyPointBox}
\end{minipage}

\IfClozeTF{\clearpage}{}
\begin{ExampleBox}{Determinante mit der Regel von Sarrus}{linear_systems_sarrus_example}
  Wir berechnen die Determinante aus
  Beispiel~\ref{ex:linear_systems_three_by_three_determinant} erneut:
  \[
    D=\det\!\begin{pmatrix}
      1&2&3\\
      4&5&6\\
      7&8&0
    \end{pmatrix}.
  \]
  \noindent\textbf{Sarrus-Schema (Nebenrechnung):}
  \[
    \begin{array}{|ccc|cc}
      1&2&3&\textcolor{green!50!black}{1}&\textcolor{green!50!black}{2}\\
      4&5&6&\textcolor{green!50!black}{4}&\textcolor{green!50!black}{5}\\
      7&8&0&\textcolor{green!50!black}{7}&\textcolor{green!50!black}{8}
    \end{array}
  \]
  \begin{ClozeArea}[32]{7}
    \[
      \begin{aligned}
        D
        &=\textcolor{THMBlue}{1\cdot5\cdot0+2\cdot6\cdot7+3\cdot4\cdot8}\\
        &\phantom{={}}\textcolor{red!80!black}{-3\cdot5\cdot7-1\cdot6\cdot8-2\cdot4\cdot0}\\[0.6em]
        &=0+84+96-105-48-0=27.
      \end{aligned}
    \]
  \end{ClozeArea}
\end{ExampleBox}

\subsection{Determinante und Lösungstyp}
\label{subsec:linear_systems_determinant_solution_criterion}

Die Cramersche Regel scheint zunächst für jedes $n\times n$-LGS eine explizite
Lösung anzugeben:
\[
  x_i=\frac{\det(A_i)}{\det(A)}
  \qquad i=1,\ldots,n.
\]
Wir wissen jedoch bereits, dass quadratische LGS auch keine oder $\infty$-viele Lösungen
besitzen können. Der gemeinsame Nenner $\det(A)$ löst diesen scheinbaren Widerspruch:

\begin{itemize}
  \item $\det(A)\neq0$: Der Quotient ist definiert und die Cramersche
        Regel liefert die eindeutige Lösung.
  \item $\det(A)=0$: Der Quotient der Cramerschen Regel ist nicht definiert und
        liefert daher auch keine eindeutige Lösung.
\end{itemize}

Damit deutet sich bereits der folgende Zusammenhang der Determinante mit dem Rang
bzw. der Lösbarkeit eines LGS an.
\begin{TheoremBox}{Determinantenkriterium für $n\times n$-LGS}{linear_systems_determinant_solution_criterion}
  Für ein $n\times n$-LGS mit Koeffizientenmatrix $A$ gilt
  \[
    \det(A)\neq0
    \quad\Longleftrightarrow\quad
    \operatorname{rang}(A)=n
    \quad\Longleftrightarrow\quad
    \exists\ \text{genau eine Lösung}.
  \]
  Durch Negation dieser Aussage erhalten wir
  \[
    \det(A)=0
    \quad\Longleftrightarrow\quad
    \operatorname{rang}(A)<n
    \quad\Longleftrightarrow\quad
    \text{keine Lösung oder } \infty\text{-viele Lösungen}.
  \]
\end{TheoremBox}

Besonders aussagekräftig ist die Determinante für homogene LGS.
\begin{KeyPointBox}{Spezialfall: homogenes $n\times n$-LGS}{linear_systems_homogeneous_determinant_criterion}
  Ein homogenes LGS besitzt immer die triviale Lösung $(0,\ldots,0)$, weshalb gilt:
  \[
    \begin{aligned}
      \det(A)\neq0
      &\quad\Longleftrightarrow\quad
      \exists\ \text{genau eine Lösung: die triviale Lösung} (0,\ldots,0),\\[0.8em]
      \det(A)=0
      &\quad\Longleftrightarrow\quad
      \infty\text{-viele Lösungen}.
    \end{aligned}
  \]
\end{KeyPointBox}
Dieses Kriterium verwenden wir später, um lineare Abhängigkeit und Unabhängigkeit von Vektoren zu prüfen.


\begin{ExampleBox}{Determinantenkriterium bei homogenen LGS}{linear_systems_homogeneous_determinant_examples}
  Wir bestimmen den Lösungstyp der folgenden homogenen LGS mithilfe der
  Determinante. Zur geometrischen Verifikation zeichnen wir auch die Geraden
  zu jeder Gleichung.

  \begin{enumerate}[label=\alph*),leftmargin=1.8em,itemsep=1.5em]
    \item
      \textbf{Eindeutige Lösung}\par
      \noindent
      \begin{minipage}[c]{0.61\linewidth}
        \vspace{-0.45\baselineskip}
        \[
          \begin{aligned}
            \textcolor{THMBlue}{x_1+x_2}&\textcolor{THMBlue}{=0},\\
            \textcolor{black!60}{x_1-x_2}&\textcolor{black!60}{=0}
          \end{aligned}
          \quad\Longleftrightarrow\quad
          [A\mid 0]=
          \left[
            \begin{array}{@{}rr|r@{}}
              1& 1&0\\
              1&-1&0
            \end{array}
          \right].
        \]
        \noindent
        Determinante der Koeffizientenmatrix:
        \[
          \det(A)=1\cdot(-1)-1\cdot1
          =\Cloze{2.2cm}{-2\neq0}
        \]
        \noindent
        $\Longrightarrow$\quad
        $\exists$ nur die triviale Lösung, d.h. $\mathcal L=\Cloze{2.8cm}{\{(0,0)\}}$
      \end{minipage}\hfill
      \begin{minipage}[c]{0.34\linewidth}
        \centering
        \CachedTikzFigure{content/resources/chapters/linear_systems/tikz/homogeneous_unique_solution_lines.tex}
      \end{minipage}

    \item
      \textbf{$\infty$-viele Lösungen}\par
      \noindent
      \begin{minipage}[c]{0.61\linewidth}
        \vspace{-0.45\baselineskip}
        \[
          \begin{aligned}
            \textcolor{THMBlue}{x_1+x_2}&\textcolor{THMBlue}{=0},\\
            \textcolor{olive!80!black}{2x_1+2x_2}&\textcolor{olive!80!black}{=0}
          \end{aligned}
          \quad\Longleftrightarrow\quad
          [A\mid 0]=
          \left[
            \begin{array}{@{}rr|r@{}}
              1&1&0\\
              2&2&0
            \end{array}
          \right].
        \]
        \noindent
        Determinante der Koeffizientenmatrix:
        \[
          \det(A)=1\cdot2-1\cdot2
          =\Cloze{2.2cm}{0}
        \]
        \noindent
        $\Longrightarrow$\quad
        \Cloze{6.2cm}{$\exists$ $\infty$-viele Lösungen.}
      \end{minipage}\hfill
      \begin{minipage}[c]{0.34\linewidth}
        \centering
        \CachedTikzFigure{content/resources/chapters/linear_systems/tikz/homogeneous_coincident_solution_lines.tex}
      \end{minipage}
  \end{enumerate}
\end{ExampleBox}


